% Propietario conservado: /Users/ruben/Documents/New project/output/EXTREMA_MEDIA_RAZON_RECIPROCIDAD_ES_EN_20260918/ARTICULO/ES/source/libro/reciprocidad_apertura.tex
\HMTDivision{\(K\;\longleftrightarrow\;\Lambda_t\;\longleftrightarrow\;\Lambda_t^*\)}{Recoverable geometry and reciprocity of coupling}{Positive deformation, lattice duality, polar structure and oriented recovery of action.}
\label{x_reciprocidad:rec:apertura}

The function of the holographic arithmetic--geometric coupling constant becomes precise by studying the transformations determined by its register. Its positional reading preserves twelve ordered components; its incidence realization distinguishes a direction; the bilateral archive allows it to be reconstructed. Composing these operations leads to a family of geometries whose fundamental volume remains constant and whose inversion produces the dual lattice. The register therefore participates both in the constitution of the geometry and in its recovery.

This development maintains the generative order established in the article. The additive and multiplicative evaluations of APP, the orientation of TRIT and the selection, transport and memory operations of TPK produce the enriched state and its joint discrete structure of the continuum. The regional coordinates and exceptional incidence are readings of that state. The vector, lattice and hyperbolic realizations that follow receive these outputs with their charts and orientations. The genealogy is preserved when the representation changes.

The distinction between scale and shape acquires an exact expression here. A deformation can dilate certain directions and contract others while preserving the total product of its factors. In the action ellipse, the product of the semiaxes preserves the action section. In the twelve radii determined by \(K\), compensation preserves the lattice covolume. In the mass operator, a decomposition separates the uniform factor from the relative transformation of determinant one. Each conservation has its own object and proof; the readers allow them to be related without confusing their domains.

Radial reciprocity articulates two complementary constructions. In the chart of the spectral double circle, inverting the radius is exactly equivalent to reflecting the complex coordinate across the line of real part \(1/2\). In the modular chart of the ellipse, exchanging semiaxes inverts the radius and changes the sign of the Klein coordinate. Both correspondences admit inverse formulas. The orientation register distinguishes the conjugate reading from a modification of the physical action section.

Substituting the full expressions for the angular readers reveals a specific arithmetic cancellation: the linear term of the preceding action section is compensated by the corresponding term of the counterangle. The regional continuation remains intact. The angular contrast, radius and hyperbolic coordinate are then expressed through the same publications involved in \(\alpha\), with the coupling constant within its dodecaphasic closure. Conversely, the oriented radial coordinate recovers the return section and allows the complete electronic expression to be rewritten.

The spectral realization in turn distinguishes two kinds of information. Lattice duality imposes a functional equation; covolume fixes the poles and their residues; deformation modifies the entire part of the completed function. The archive recovers the direction determining that variation and provides explicit bounds for the reconstruction error. Conservation and reversibility are thus expressed through operations encompassing the register, geometry and their analytic readings.

\HMTMathematicalPaper
